\chapter{Structuurbepaling: \texorpdfstring{$^{13}$C}{13C}-NMR, MS en IR samen}\label{ch:b2:structure-determination}

Een geurstoffenlaboratorium krijgt een flesje met een kleurloze vloeistof die naar jasmijn ruikt. Niemand heeft
er een etiket op geschreven. Een \omterm{def:b2:mass-spec-atomic:spectrum}{massaspectrum} geeft de molecuulformule, een infraroodspectrum de functionele
groepen, een koolstof-13-spectrum het skelet en een protonspectrum de manier waarop de stukken verbonden zijn; binnen een
uur heeft de vloeistof een naam. Dit hoofdstuk voegt \omterm{def:b2:structure-determination:c13}{koolstof-13-NMR} toe aan de proton-NMR en de infraroodspectroscopie
van het deel van jaar~1 en aan de \omterm{def:b2:mass-spec-atomic:spectrum}{massaspectrometrie} van \cref{ch:b2:mass-spec-atomic}, en maakt van de vier
één methode.

\begin{recall}
Het deel van jaar~1: chemische verschuiving en afscherming, equivalente protonen, integratie, spin-spinkoppeling
en de $n + 1$-regel, infraroodgolfgetallen van functionele groepen en het vingerafdrukgebied, de
onverzadigingsgraad, een structuur oplossen uit drie spectra. \Cref{ch:b2:mass-spec-atomic}: het
\omterm{def:b2:mass-spec-atomic:molecular-ion}{molecuulion}, \omterm{def:b2:mass-spec-atomic:isotope-pattern}{isotopenpatronen}, fragmentatie.
\end{recall}

\section{Koolstof-13-NMR}

Koolstof-12, het hoofdisotoop, heeft geen kernspin en geeft geen NMR-signaal. Koolstof-13 heeft een spin van
$\tfrac12$, zoals het proton, maar het vormt maar 1.1~\% van natuurlijke koolstof en zijn signaal is, kern voor kern,
zwakker dan dat van een proton; een koolstofspectrum vraagt daarom meer monster of meer scans.
% ledger: iso:C13

\begin{definition}[Koolstof-13-NMR]\label{def:b2:structure-determination:c13}
\emph{Koolstof-13-NMR}\index{koolstof-13-NMR} registreert de resonanties van de \ce{^{13}C}-kernen van een monster.
Meestal gebeurt dat met \emph{breedbandontkoppeling}\index{breedbandontkoppeling}: de protonen worden
tijdens de meting over hun hele frequentiebereik bestraald, wat elke
koolstof-protonkoppeling wegneemt, zodat elke soort koolstofatoom één enkele lijn geeft. \emph{Equivalente
koolstofatomen}\index{equivalente koolstofatomen}, die door een symmetrie van het molecuul (of door een snelle rotatie) in elkaar overgaan,
geven dezelfde lijn.
\end{definition}

\begin{proposition}[Geen koolstof-koolstofkoppeling]\label{prop:b2:structure-determination:no-cc-coupling}
Koppelingen tussen naburige koolstofatomen verschijnen niet in een \ce{^{13}C}-spectrum bij natuurlijke
abundantie.
\end{proposition}

\begin{proof}
Een koppeling tussen twee koolstofatomen is alleen zichtbaar in moleculen waarin beide \ce{^{13}C} zijn. Voor twee gegeven
naburige posities is de kans $a_{13}^2 = 0.011^2 \approx 1.2 \times 10^{-4}$: ongeveer één
molecuul op achtduizend. Het signaal van elk koolstofatoom komt bijna volledig van moleculen waarin zijn
buren \ce{^{12}C} zijn, onzichtbaar voor de spectrometer; de gekoppelde satellieten zijn honderd keer
zwakker dan de ruis van een gewoon spectrum.
\end{proof}

\begin{proposition}[Aantal lijnen]\label{prop:b2:structure-determination:signals}
Een ontkoppeld \ce{^{13}C}-spectrum toont één lijn per set \omterm{def:b2:structure-determination:c13}{equivalente koolstofatomen}, tenzij twee sets toevallig
dezelfde verschuiving hebben.
\end{proposition}

\begin{proof}[Redenering]
\omterm{def:b2:structure-determination:c13}{Equivalente koolstofatomen} bevinden zich in identieke omgevingen, dus ze hebben dezelfde verschuiving; niet-equivalente koolstofatomen
verschillen meestal. Zonder koppelingen met protonen en zonder koppelingen tussen koolstofatomen
(\cref{prop:b2:structure-determination:no-cc-coupling}) splitst niets een lijn. Toevallige overlapping
treedt op wanneer twee verschillende koolstofatomen verschuivingen hebben die dichter bij elkaar liggen dan de resolutie, wat vaak gebeurt bij de
koolstofatomen van een benzeenring rond \qty{128}{ppm}.
\end{proof}

De verschuivingen spreiden zich uit over ongeveer \qty{220}{ppm}, twintig keer het bereik van protonen, dus overlappingen zijn zeldzamer dan
in protonspectra, en de verschuiving alleen zegt al welke soort koolstofatoom het is. De grafiek hieronder toont de verschuivingen
die gemeten zijn voor de verbindingen van dit hoofdstuk.

\begin{omfigure}
\begin{tikzpicture}
\begin{axis}[width=0.86\linewidth, height=4.6cm, xmin=0, xmax=220, ymin=-0.6, ymax=4.6, x dir=reverse,
  xlabel={$\delta$ (ppm)}, label style={font=\small}, tick label style={font=\scriptsize},
  ytick={0,1,2,3,4}, yticklabels={alkyl-C, C--O, aromatisch C, ester-C=O, keton-C=O},
  xtick={0,20,40,60,80,100,120,140,160,180,200,220}, grid=major, grid style={black!10}]
\addplot[only marks, mark=|, mark size=5pt, thick, omDef] table[x=shift, y=cls] {figdata/out/bachelor-2-nmr-spectra-d.dat};
\end{axis}
\end{tikzpicture}

{\small Gemeten \ce{^{13}C}-verschuivingen van butaan-2-on, pentaan-2-on, ethylbenzoaat en benzylethanoaat,
gesorteerd naar soort koolstofatoom (CDCl$_3$; gegevens van PubChem). Ketoncarbonylen liggen rond 209, estercarbonylen rond
167--171, \omterm{def:b2:huckel:aromatic}{aromatische} koolstofatomen tussen 128 en 137, koolstofatomen gebonden aan een esterzuurstofatoom rond 61--66, en
alkylkoolstofatomen onder 46.}
% figdata: bachelor-2/nmr-spectra.py part d
% ledger: c13:butanone, c13:pentan-2-one, c13:ethylbenzoate, c13:benzylacetate
\end{omfigure}

\begin{proposition}[Hoogten zijn geen aantallen]\label{prop:b2:structure-determination:intensities}
In een routinematig ontkoppeld \ce{^{13}C}-spectrum zijn de hoogten van de lijnen niet evenredig met de
aantallen koolstofatomen die ze voorstellen; koolstofatomen zonder waterstofatomen geven zwakke lijnen.
\end{proposition}

\begin{proof}[Redenering]
Routinespectra herhalen de meting sneller dan de traagste koolstofatomen naar het evenwicht terugrelaxeren, en
koolstofatomen zonder gebonden protonen relaxeren traag, zodat hun signaal gedeeltelijk verzadigd is. Ontkoppeling
versterkt ook de signalen van koolstofatomen die protonen dragen, met een hoeveelheid die van hun omgeving afhangt.
In butaan-2-on is de carbonyllijn de zwakste van de vier, al staat ze net als de andere voor één koolstofatoom.
Integratie, zo nuttig voor protonen, wordt daarom niet gebruikt voor routinematige koolstofspectra.
\end{proof}

\section{DEPT}

\begin{definition}[DEPT]\label{def:b2:structure-determination:dept}
\emph{DEPT}\index{DEPT} (distortionless enhancement by polarisation transfer) is een set koolstof-13-experimenten
die de signalen van koolstofatomen met gebonden protonen een teken geven of laten verdwijnen, naargelang
hun aantal waterstofatomen. Een \emph{quaternair koolstofatoom}\index{quaternair koolstofatoom} draagt geen waterstof
(gebonden aan vier andere atomen, of een carbonyl- of \omterm{def:b2:huckel:aromatic}{aromatisch} koolstofatoom zonder H); het geeft geen DEPT-signaal.
\end{definition}

\begin{proposition}[DEPT lezen]\label{prop:b2:structure-determination:dept}
Bij DEPT-135 wijzen \ce{CH}- en \ce{CH3}-koolstofatomen omhoog, \ce{CH2}-koolstofatomen omlaag en zijn quaternaire koolstofatomen
afwezig. Bij DEPT-90 verschijnen alleen \ce{CH}-koolstofatomen.
\end{proposition}

\admitted

Het experiment draagt magnetisatie over van de protonen naar het koolstofatoom waaraan ze gebonden zijn, via een
reeks radiofrequente pulsen waarvan de laatste hoek de respons selecteert; hoe de pulsen dat doen, wordt
uitgelegd in het deel van jaar~3. Het ontkoppelde spectrum vergelijken met DEPT-135 en DEPT-90 deelt elke
lijn in als C, CH, \ce{CH2} of \ce{CH3}.

\begin{omfigure}
\pgfplotsset{dept/.style={width=\linewidth, height=2.9cm, ycomb, x dir=reverse, ymin=-1, ymax=1.15,
  ytick=\empty, axis y line=none, axis x line=middle, x axis line style={-}, tick label style={font=\tiny},
  title style={font=\scriptsize, at={(0.0,0.95)}, anchor=north west}, every axis plot/.append style={thick}}}
\begin{minipage}[t]{0.49\linewidth}\centering
{\small butaan-2-on}\par
\begin{tikzpicture}
\begin{axis}[dept, xmin=0, xmax=220, xtick={0,50,100,150,200}, title={ontkoppeld}]
\addplot[omDef] table[x=shift, y=dec] {figdata/out/bachelor-2-nmr-spectra-a.dat};
\end{axis}
\end{tikzpicture}\par
\begin{tikzpicture}
\begin{axis}[dept, xmin=0, xmax=220, xtick={0,50,100,150,200}, title={DEPT-135}]
\addplot[omThm] table[x=shift, y=d135] {figdata/out/bachelor-2-nmr-spectra-a.dat};
\end{axis}
\end{tikzpicture}\par
\begin{tikzpicture}
\begin{axis}[dept, xmin=0, xmax=220, xtick={0,50,100,150,200}, title={DEPT-90}]
\addplot[omMeth] table[x=shift, y=d90] {figdata/out/bachelor-2-nmr-spectra-a.dat};
\end{axis}
\end{tikzpicture}
\end{minipage}\hfill
\begin{minipage}[t]{0.49\linewidth}\centering
{\small ethylbenzoaat}\par
\begin{tikzpicture}
\begin{axis}[dept, xmin=0, xmax=220, xtick={0,50,100,150,200}, title={ontkoppeld}]
\addplot[omDef] table[x=shift, y=dec] {figdata/out/bachelor-2-nmr-spectra-b.dat};
\end{axis}
\end{tikzpicture}\par
\begin{tikzpicture}
\begin{axis}[dept, xmin=0, xmax=220, xtick={0,50,100,150,200}, title={DEPT-135}]
\addplot[omThm] table[x=shift, y=d135] {figdata/out/bachelor-2-nmr-spectra-b.dat};
\end{axis}
\end{tikzpicture}\par
\begin{tikzpicture}
\begin{axis}[dept, xmin=0, xmax=220, xtick={0,50,100,150,200}, title={DEPT-90}]
\addplot[omMeth] table[x=shift, y=d90] {figdata/out/bachelor-2-nmr-spectra-b.dat};
\end{axis}
\end{tikzpicture}
\end{minipage}

{\small Ontkoppelde \ce{^{13}C}-spectra (gemeten verschuivingen en hoogten, gegevens van PubChem) met de
DEPT-responsen (tekens volgens \cref{prop:b2:structure-determination:dept}, hoogten schematisch). Butaan-2-on: C=O
209.3 (afwezig bij \omterm{def:b2:structure-determination:dept}{DEPT}), \ce{CH2} 36.9 (omlaag), \ce{CH3} 29.4 en 7.9 (omhoog). Ethylbenzoaat: C=O 166.5 en
het ringkoolstofatoom dat de ester draagt, 130.6, verdwijnen; drie \omterm{def:b2:huckel:aromatic}{aromatische} CH (132.8, 129.6, 128.3) blijven bij
DEPT-90; \ce{OCH2} 60.9 wijst omlaag.}
% figdata: bachelor-2/nmr-spectra.py parts a, b
\end{omfigure}

\section{De technieken combineren}

Elke techniek beantwoordt haar eigen vraag, en de antwoorden worden in een vaste volgorde samengevoegd: eerst de formule,
omdat die al de rest begrenst; dan de functionele groepen en het skelet; dan de
verbindingen.

\begin{method}[Een onbekende oplossen]\label{met:b2:structure-determination:solve}
\begin{enumerate}
  \item Formule: \omterm{def:b2:mass-spec-atomic:molecular-ion}{molecuulion}, $M+1$ voor het aantal koolstofatomen, \omterm{def:b2:mass-spec-atomic:isotope-pattern}{isotopenpatronen}, \omterm{def:b2:mass-spec-atomic:nitrogen-rule}{stikstofregel}; een
        exacte massa als die beschikbaar is. Onverzadigingsgraad.
  \item Infrarood: \ce{C=O} (en welke soort), \ce{O-H}, \ce{N-H}, \ce{C#N}, \omterm{def:b2:huckel:aromatic}{aromatische} of alkeen-\ce{C-H}
        boven \qty{3000}{cm^{-1}}.
  \item Koolstof-13 en \omterm{def:b2:structure-determination:dept}{DEPT}: aantal soorten koolstofatomen (symmetrie), hun klassen uit de verschuivingen, en
        C, CH, \ce{CH2}, \ce{CH3} uit \omterm{def:b2:structure-determination:dept}{DEPT}. Vergelijk met de formule.
  \item Proton-NMR: integraties, verschuivingen, multipliciteiten; buren uit de $n + 1$-regel.
  \item Voeg de fragmenten samen; schrijf elke kandidaatstructuur op.
  \item Toets elke kandidaat aan elk gegeven, ook aan de fragmenten van het \omterm{def:b2:mass-spec-atomic:spectrum}{massaspectrum}; behoud degene
        die ze allemaal verklaart.
\end{enumerate}
\end{method}

\begin{omfigure}
\begin{tikzpicture}[x=1cm, y=1cm, st/.style={draw=black!70, rounded corners=2pt, fill=white,
  font=\small, align=center, minimum height=0.9cm, inner sep=3pt, text width=2.3cm}]
\node[st, fill=omDef!10] (ms) at (0,0) {MS\\formule, $M+1$};
\node[st] (dbe) at (2.9,0) {onverzadigings-\\graad};
\node[st, fill=omDef!10] (ir) at (5.8,0) {IR\\functionele groepen};
\node[st, fill=omDef!10] (c13) at (8.7,0) {\ce{^{13}C}, DEPT\\skelet};
\node[st, fill=omDef!10] (h1) at (8.7,-1.8) {\ce{^1H}\\verbindingen};
\node[st] (cand) at (5.8,-1.8) {kandidaat-\\structuren};
\node[st, fill=omThm!10] (chk) at (2.9,-1.8) {elk gegeven\\toetsen};
\node[st] (ans) at (0,-1.8) {structuur};
\draw[->, thick] (ms) -- (dbe); \draw[->, thick] (dbe) -- (ir); \draw[->, thick] (ir) -- (c13);
\draw[->, thick] (c13) -- (h1); \draw[->, thick] (h1) -- (cand); \draw[->, thick] (cand) -- (chk);
\draw[->, thick] (chk) -- (ans);
\draw[->, thick, densely dashed] (chk.south) -- ++(0,-0.5) -| node[pos=0.25, below, font=\footnotesize]
  {een onverklaard gegeven: terug naar de kandidaten} (cand.south);
\end{tikzpicture}
\par\vspace{4pt}

{\small De werkvolgorde voor een onbekende (\cref{met:b2:structure-determination:solve}). De laatste stap
wordt het vaakst overgeslagen: een structuur die één piek onverklaard laat, is nog niet het antwoord.}
\end{omfigure}

\section{Uitgewerkte gevallen}

\begin{example}[Een keton, C$_5$H$_{10}$O]\label{ex:b2:structure-determination:ketone}
Een onbekende heeft $M = 86$, een IR-band rond \qty{1715}{cm^{-1}}, koolstoflijnen bij 208.9, 45.7, 29.8, 17.4 en
13.7~ppm, en protonsignalen bij 2.41 (t, 2H), 2.13 (s, 3H), 1.61 (sextet, 2H) en 0.91 (t, 3H).
Eén onverzadiging, een verzadigd keton (geen aldehydeproton rond 9.7). Vijf koolstofatomen, vijf lijnen:
geen symmetrie. Het singlet van drie protonen bij 2.13 is een \ce{CH3} naast de carbonylgroep; de opeenvolging triplet,
sextet, triplet is een propylgroep waarvan de \ce{CH2} bij 2.41 aan de carbonylgroep grenst. Pentaan-2-on,
\ce{CH3COCH2CH2CH3}. Zijn symmetrische isomeer pentaan-3-on zou drie koolstoflijnen tonen en geen singlet;
zijn \omterm{def:b2:mass-spec-atomic:spectrum}{massaspectrum} heeft het McLafferty-ion bij 58 dat pentaan-3-on mist (\cref{ch:b2:mass-spec-atomic}).
% ledger: c13:pentan-2-one, nmr:pentan-2-one, ir:CO-ketone, ms:pentan-2-one
\end{example}

\begin{example}[Een aromatische ester, C$_9$H$_{10}$O$_2$]\label{ex:b2:structure-determination:ester}
Een onbekende met formule \ce{C9H10O2} (onverzadigingsgraad 5) toont een carbonylband en zeven
koolstoflijnen: 166.5 (C), 132.8 (CH), 130.6 (C), 129.6 (CH), 128.3 (CH), 60.9 (\ce{CH2}) en 14.3
(\ce{CH3}). Zijn protonspectrum: 8.05 (m, 2H), 7.35--7.63 (m, 3H), 4.37 (q, 2H), 1.38 (t, 3H). Een
monogesubstitueerde benzeenring (vier \omterm{def:b2:huckel:aromatic}{aromatische} koolstoflijnen, waarvan twee voor elk twee koolstofatomen; vijf
\omterm{def:b2:huckel:aromatic}{aromatische} protonen) verklaart vier onverzadigingen, de carbonylgroep de vijfde. Het paar kwartet--triplet
is een ethylgroep; zijn \ce{CH2} bij 4.37 (en bij 60.9 in het koolstofspectrum) is aan zuurstof gebonden. De estercarbonyl
bij 166.5 is aan de ring gebonden: de twee protonen bij 8.05, ernaast, worden naar lager veld geduwd.
Ethylbenzoaat, \ce{C6H5COOCH2CH3}.
% ledger: c13:ethylbenzoate, nmr:ethylbenzoate
\end{example}

\section{Valkuilen}

\begin{itemize}
  \item \emph{Uitwisselbare protonen} (\ce{O-H}, \ce{N-H}) geven brede signalen bij variabele verschuivingen, zijn
        vaak niet gekoppeld aan hun buren, en verdwijnen wanneer een druppel \ce{D2O} met het
        monster geschud wordt.
  \item \emph{Overlappende signalen}: twee koolstofatomen van een benzeenring, of verscheidene \ce{CH2} van een keten, kunnen
        één lijn delen; tel de lijnen tegenover de formule voor je over symmetrie besluit.
  \item \emph{Diastereotope protonen}: de twee protonen van een \ce{CH2} naast een stereocentrum zijn niet
        equivalent, kunnen verschillende verschuivingen hebben en met elkaar koppelen.
  \item \emph{Tweede-ordepatronen}: wanneer twee gekoppelde protonen verschuivingen hebben die minder dan enkele keren hun
        koppelingsconstante van elkaar liggen, buigen de multipletten naar elkaar toe en faalt de $n+1$-regel; een
        spectrometer met hoger veld herstelt eenvoudige patronen.
\end{itemize}

\begin{safety}{\ghs[5mm]{GHS02}\,\ghs[5mm]{GHS07}\,\ghs[5mm]{GHS08}}
Butaan-2-on is ontvlambaar. De magneet van een NMR-spectrometer staat altijd aan: geen stalen gereedschap,
gasflessen of magnetische kaarten in de buurt, en geen toegang voor mensen met een pacemaker.
% ledger: ghs:butanone, ghs:benzylacetate
\end{safety}

\section{Oefeningen}

\begin{exercise}[$\star$]\label{exo:b2:structure-determination:1}
Hoeveel lijnen toont het ontkoppelde \ce{^{13}C}-spectrum van elk xyleen (1,2-, 1,3- en
1,4-dimethylbenzeen)?
\end{exercise}

\begin{exercise}[$\star$]\label{exo:b2:structure-determination:2}
Voorspel de DEPT-135- en DEPT-90-spectra van butaan-2-ol.
\end{exercise}

\begin{exercise}[$\star$]\label{exo:b2:structure-determination:3}
Bij welke soort koolstofatoom hoort een lijn bij 209, 170, 129, 64 en 14~ppm het waarschijnlijkst?
\end{exercise}

\begin{exercise}[$\star$]\label{exo:b2:structure-determination:4}
Bereken de onverzadigingsgraad van \ce{C8H8O2} en stel een structuur voor die een benzeenring bevat.
\end{exercise}

\begin{exercise}[$\star\star$]\label{exo:b2:structure-determination:5}
Een vloeistof \ce{C4H8O} toont een IR-band bij \qty{1715}{cm^{-1}} en koolstoflijnen bij 209, 37, 29 en 8~ppm
(DEPT-135: 37 omlaag, 29 en 8 omhoog, 209 afwezig). Identificeer ze.
\end{exercise}

\begin{exercise}[$\star\star$]\label{exo:b2:structure-determination:6}
Hoe zou je pentaan-2-on van pentaan-3-on onderscheiden met alleen \ce{^{13}C}-NMR? Met proton-NMR? Met
\omterm{def:b2:mass-spec-atomic:spectrum}{massaspectrometrie}?
\end{exercise}

\begin{exercise}[$\star\star$]\label{exo:b2:structure-determination:7}
Een ester heeft $M = 88$ en protonsignalen bij 4.12 (q, 2H), 2.04 (s, 3H) en 1.26 (t, 3H) (gegevens voor de
oefening). Identificeer hem.
\end{exercise}

\begin{exercise}[$\star\star$]\label{exo:b2:structure-determination:8}
Onderscheid 1,2-dichloorbenzeen van 1,4-dichloorbenzeen met \ce{^{13}C}-NMR.
\end{exercise}

\begin{exercise}[$\star\star$]\label{exo:b2:structure-determination:9}
Waarom zijn in 2-chloorbutaan de twee protonen van de groep \ce{CH2} niet equivalent? Wat doet dat met
het protonspectrum?
\end{exercise}

\begin{exercise}[$\star\star\star$]\label{exo:b2:structure-determination:10}
Een onbekende \ce{C9H10O} toont een IR-band bij \qty{1690}{cm^{-1}}, koolstoflijnen bij 200.8 (C), 137.0 (C),
132.9 (CH), 128.6 (CH), 128.0 (CH), 31.8 (\ce{CH2}) en 8.3 (\ce{CH3}), en protonsignalen bij 7.95 (m,
2H), 7.40--7.55 (m, 3H), 2.98 (q, 2H) en 1.22 (t, 3H) (gegevens voor de oefening). Identificeer ze, en verklaar het
lage golfgetal van de carbonylgroep.
\end{exercise}

\begin{exercise}[$\star\star\star$]\label{exo:b2:structure-determination:11}
Vier isomeren \ce{C5H10O2}: pentaanzuur, methylbutanoaat, ethylpropanoaat en propylethanoaat.
Geef voor elk een kenmerk van zijn IR- of protonspectrum dat het van de andere drie onderscheidt.
\end{exercise}

\begin{exercise}[$\star\star\star$]\label{exo:b2:structure-determination:12}
Een rapport kent de lijn bij 166.5~ppm van ethylbenzoaat toe aan het ringkoolstofatoom dat de ester draagt, en
de lijn bij 130.6~ppm aan het carbonylkoolstofatoom. Toon met de DEPT-gegevens en de verschuivingsgrafiek aan dat die
toekenning niet kan kloppen.
\end{exercise}

\section{Probleem: vier spectra van jasmijn}

\begin{problem}[{Weekendprobleem --- de formule uit het molecuulion, de carbonylband in het infrarood,
koolstof-13 en DEPT, het protonspectrum, en een structuur die aan de fragmenten getoetst wordt}]
\label{pb:b2:structure-determination:1}
Een kleurloze vloeistof met een bloemige jasmijngeur geeft de volgende gegevens. \omterm{def:b2:mass-spec-atomic:spectrum}{Massaspectrum} (NIST
WebBook): $m/z$ 150 (31.7), 151 (3.1), 108 (100), 91 (71.7), 90 (40.6), 43 (37.6). Infrarood (gegevens voor de
oefening): sterke banden bij 1740 en \qty{1230}{cm^{-1}}, zwakke banden net boven \qty{3000}{cm^{-1}}, geen
brede band rond \qty{3300}{cm^{-1}}. Koolstof-13, CDCl$_3$ (gegevens van PubChem): 170.70, 136.14, 128.56,
128.24, 66.24, 20.82~ppm. Proton, CDCl$_3$, \qty{90}{MHz}: 7.33 (m, 5H), 5.09 (s, 2H), 2.06 (s, 3H).
% ledger: use:benzylacetate, ms:benzylacetate, c13:benzylacetate, nmr:benzylacetate, ir:CO-ester, iso:C12, iso:C13, mass:C12, mass:H1, mass:O16

\begin{center}
\begin{tikzpicture}
\pgfplotsset{dept/.style={width=0.8\linewidth, height=3.2cm, ycomb, x dir=reverse, ymin=0, ymax=1.15,
  xmin=0, xmax=200, ytick=\empty, axis y line=none, axis x line=bottom, x axis line style={-}, tick label style={font=\tiny},
  title style={font=\scriptsize, at={(0.02,0.82)}, anchor=west}, every axis plot/.append style={thick}}}
\begin{axis}[dept, title={ontkoppeld}, xtick={0,20,40,60,80,100,120,140,160,180,200}]
\addplot[omDef] table[x=shift, y=dec] {figdata/out/bachelor-2-nmr-spectra-c.dat};
\end{axis}
\end{tikzpicture}
\end{center}
% figdata: bachelor-2/nmr-spectra.py part c

\medskip
\noindent\textbf{Deel I --- De formule.}
\begin{enumerate}
  \item Schat uit de verhouding 151/150 het aantal koolstofatomen.
  \item Wat zegt de \omterm{def:b2:mass-spec-atomic:nitrogen-rule}{stikstofregel}?
  \item Wat blijft er met negen koolstofatomen over van de massa 150? Stel een formule met zuurstof voor.
  \item Waarom is \ce{C10H14O}, ook met nominale massa 150, minder waarschijnlijk?
  \item Bereken de onverzadigingsgraad.
  \item Bereken de \omterm{def:b2:mass-spec-atomic:exact-mass}{monoisotopische massa} van de weerhouden formule.
\end{enumerate}

\medskip
\noindent\textbf{Deel II --- Infrarood.}
\begin{enumerate}[resume]
  \item Wijs de band bij \qty{1740}{cm^{-1}} toe. Op welke familie wijst zijn positie?
  \item Wijs de band bij \qty{1230}{cm^{-1}} toe.
  \item Waarop wijzen de zwakke banden boven \qty{3000}{cm^{-1}}?
  \item Wat sluit de afwezigheid van een brede band rond \qty{3300}{cm^{-1}} uit?
  \item Is de carbonylgroep geconjugeerd met de ring? Leg uit vanuit haar golfgetal.
\end{enumerate}

\medskip
\noindent\textbf{Deel III --- Koolstof-13 en \omterm{def:b2:structure-determination:dept}{DEPT}.}
\begin{enumerate}[resume]
  \item Wijs elke lijn toe aan een soort koolstofatoom.
  \item Welke lijn wijst omlaag bij DEPT-135?
  \item Welke lijnen verdwijnen bij DEPT-135?
  \item Welke lijnen blijven over bij DEPT-90?
  \item Hoeveel soorten koolstofatomen heeft een monogesubstitueerde benzeenring? Leg uit waarom er maar twee \omterm{def:b2:huckel:aromatic}{aromatische}
        CH-lijnen verschijnen.
  \item De lijn bij 128.24 is de hoogste. Staat ze voor de meeste koolstofatomen?
\end{enumerate}

\medskip
\noindent\textbf{Deel IV --- Proton-NMR en structuur.}
\begin{enumerate}[resume]
  \item Wijs de drie protonsignalen toe.
  \item Waarom zijn de signalen bij 5.09 en 2.06 singlets?
  \item Waarom ligt de \ce{CH2} bij 5.09 zo ver naar lager veld?
  \item Stel de structuur samen en geef de naam.
  \item Zijn isomeer methyl-2-fenylethanoaat, \ce{C6H5CH2COOCH3}, heeft dezelfde formule. Welk gegeven sluit het
        uit?
  \item Verklaar de fragmenten bij 108, 91 en 43.
  \item Hoeveel lijnen zou een perfect opgelost ontkoppeld \ce{^{13}C}-spectrum van de verbinding tonen, en
        hoeveel ervan wijzen omlaag bij DEPT-135?
\end{enumerate}
\end{problem}
